\chapter{Sets, lines, spaces, structures, fields, and rings.}

\section{Sets.}

A set is simply a collection of different numbers. For example, the set of all integers which are greater than or equal
to 1 and less than or equal to 5, could simply be written as;

\begin{equation*}
    \left\{1,2,3,4,5\right\}
\end{equation*}
\vspace{0\baselineskip}

Alternatively, we could define this set using mathematical set-builder notation, by writing it as follows;

\begin{equation*}
    \left\{ \, x \in \mathbb{Z}\;\, \mid \;\, 1 \leq x \leq 6 \, \right\}
\end{equation*}
\vspace{0\baselineskip}

Now consider all of the real numbers. They can be thought of as belonging to or forming a set of numbers which
mathematicians denote as \(\mathbb{R}\). Using mathematical set-builder notation, we can define this set as follows;

\begin{equation*}
    \mathbb{R} = \left\{ \, x \;\, | \;\, x \in \mathbb{R} \, \right\}
\end{equation*}


\section{Lines.}

Certain sets of numbers can be imagined as living or residing on a line. When it comes to imagining the set of real
numbers \(\mathbb{R}\) on a line, mathematicians call this line the real number line.


\section{Spaces.}

Now consider what would happen if we took two of these real number lines and arranged them such that they;

\begin{itemize}[itemsep=0pt] % [label=\alph*)]
    \item were perpendicular to each other, and
    \item both intersected each other at 0 on their axis.
\end{itemize}

If we did this, then we would end up with a 2-dimensional real number space; that is, a 2-dimensional real number plane,
or a real Euclidean plane. In such a case, the two axes which comprise the plane would usually be denoted as \(x\) and
\(y\), while the resulting plane is denoted as \(\mathbb{R}^2\). \\

Using mathematical set-builder notation, we can define \(\mathbb{R}^2\) as follows;

\begin{equation*}
    \mathbb{R}^2 = \left\{ \, (x,y) \;\, | \;\, x,y \in \mathbb{R} \, \right\}
\end{equation*}
\vspace{0\baselineskip}

The trouble with \(\mathbb{R}^2\), is that it doesn't really do much other than define a set of pairs of values, all of
which are of the form \((x,y)\). For example, it doesn't define any operations which allow us to operate on any of the
values within the set that it defines. What if we wanted to access one its values or add two of its values together.
How would we do this? This is where structures come in.


\section{Structures.}

In contemporary mathematics, a structure -- or more precisely an algebraic structure or a mathematical structure, is
defined as;

\begin{framed}
\begin{quote}
    a set together with a collection of operations -- and sometimes relations or distinguished constants, that satisfy a
    list of axioms.
\end{quote}
    \end{framed}


\section{Fields.}



\section{Rings.}